\section{\texorpdfstring{\(R_{e/\mu}\) Correction Factors}{Re/mu Correction Factors}}
\label{sec:remu-corrections}

This appendix collects the correction conventions used in
\secref{sec:measurement-strategy}. The purpose is to separate the experimental
strategy from the algebraic bookkeeping, since additional small corrections can
be added here as the analysis matures.

\subsection{High- and Low-Energy Counts}
\label{subsec:remu-high-low-counts}

Define the reconstructed high- and low-energy samples by an energy cut
\(E_{\mathrm{cut}}\),
\begin{align}
  N_H &\equiv N(E_{\mathrm{reco}}>E_{\mathrm{cut}}), &
  N_L &\equiv N(E_{\mathrm{reco}}<E_{\mathrm{cut}}).
\end{align}

Let \(N_H^{\pi e}\) and \(N_L^{\pi e}\) denote true
\(\pi\to e\nu\) events reconstructed above and below threshold. Similarly,
let \(N_H^{\pi\mu}\) and \(N_L^{\pi\mu}\) denote
\(\pi\to\mu\to e\) events reconstructed above and below threshold. The observed
high-to-low ratio is therefore
\begin{equation}
  \frac{N_H}{N_L}
  =
  \frac{N_H^{\pi e}+N_H^{\pi\mu}}
       {N_L^{\pi e}+N_L^{\pi\mu}} .
  \label{eq:observed-high-low-ratio}
\end{equation}

The true branching-ratio ratio before detector effects is
\begin{equation}
  R_{e/\mu}^{\mathrm{true}}
  =
  \frac{N_H^{\pi e}+N_L^{\pi e}}
       {N_H^{\pi\mu}+N_L^{\pi\mu}} .
  \label{eq:true-remu-counts}
\end{equation}

The high-energy threshold is chosen so that ordinary stopped-muon Michel events
above threshold are negligible at the required precision,
\begin{equation}
  N_H^{\pi\mu} \simeq 0 .
  \label{eq:negligible-muon-leakage}
\end{equation}

With this approximation,
\eqnref{eq:observed-high-low-ratio} becomes
\begin{equation}
  \frac{N_H}{N_L}
  \simeq
  \frac{N_H^{\pi e}}
       {N_L^{\pi e}+N_L^{\pi\mu}},
  \label{eq:observed-ratio-tail-approx}
\end{equation}
while \eqnref{eq:true-remu-counts} becomes
\begin{equation}
  R_{e/\mu}^{\mathrm{true}}
  \simeq
  \frac{N_H^{\pi e}+N_L^{\pi e}}
       {N_L^{\pi\mu}} .
  \label{eq:true-ratio-tail-approx}
\end{equation}

The leading difference between these expressions is the number of true
\(\pi\to e\nu\) events reconstructed below the energy threshold.

\subsection{Tail Correction}
\label{subsec:remu-tail-correction}

The observed high-to-low ratio is
\begin{equation}
  \frac{N_H}{N_L}
  =
  \frac{N_H^{\pi e}+N_H^{\pi\mu}}
       {N_L^{\pi e}+N_L^{\pi\mu}} .
\end{equation}

The threshold is chosen so that ordinary Michel events reconstructed above the
cut are negligible at the required precision,
\begin{equation}
  N_H^{\pi\mu}\simeq0 ,
\end{equation}
giving
\begin{equation}
  \frac{N_H}{N_L}
  \simeq
  \frac{N_H^{\pi e}}
       {N_L^{\pi e}+N_L^{\pi\mu}} .
  \label{eq:observed-tail-ratio}
\end{equation}

The true branching-ratio ratio is
\begin{equation}
  R_{e/\mu}^{\mathrm{true}}
  =
  \frac{N_H^{\pi e}+N_L^{\pi e}}
       {N_H^{\pi\mu}+N_L^{\pi\mu}}
  \simeq
  \frac{N_H^{\pi e}+N_L^{\pi e}}
       {N_L^{\pi\mu}} .
\end{equation}

Insert
\begin{equation}
  1
  =
  \frac{N_H/N_L}{N_H/N_L}
\end{equation}
to obtain
\begin{align}
  R_{e/\mu}^{\mathrm{true}}
  &=
  \left[
    \frac{N_H}{N_L}
    \cdot
    \frac{N_L}{N_H}
  \right]
  \cdot
  \frac{N_H^{\pi e}+N_L^{\pi e}}
       {N_L^{\pi\mu}}
  \\
  &=
  \frac{N_H}{N_L}
  \left[
    \frac{N_L}{N_H}
    \cdot
    \frac{N_H^{\pi e}+N_L^{\pi e}}
         {N_L^{\pi\mu}}
  \right].
\end{align}

Using
\eqnref{eq:observed-tail-ratio},
\begin{equation}
  \frac{N_L}{N_H}
  \simeq
  \frac{N_L^{\pi e}+N_L^{\pi\mu}}
       {N_H^{\pi e}},
\end{equation}
so
\begin{align}
  R_{e/\mu}^{\mathrm{true}}
  &\simeq
  \frac{N_H}{N_L}
  \left[
    \frac{
      (N_L^{\pi e}+N_L^{\pi\mu})
      (N_H^{\pi e}+N_L^{\pi e})
    }
    {
      N_H^{\pi e}N_L^{\pi\mu}
    }
  \right]
  \\
  &=
  \frac{N_H}{N_L}
  \left[
    \frac{
      N_L^{\pi e}N_H^{\pi e}
      +
      (N_L^{\pi e})^2
      +
      N_L^{\pi\mu}N_H^{\pi e}
      +
      N_L^{\pi\mu}N_L^{\pi e}
    }
    {
      N_H^{\pi e}N_L^{\pi\mu}
    }
  \right].
\end{align}

Dividing each term separately,
\begin{align}
  R_{e/\mu}^{\mathrm{true}}
  &=
  \frac{N_H}{N_L}
  \Bigg[
    \frac{N_L^{\pi e}}
         {N_L^{\pi\mu}}
    +
    \frac{(N_L^{\pi e})^2}
         {N_H^{\pi e}N_L^{\pi\mu}}
    +
    1
    +
    \frac{N_L^{\pi e}}
         {N_H^{\pi e}}
  \Bigg].
\end{align}

The first two correction terms are suppressed by additional factors of
\(R_{e/\mu}\sim10^{-4}\),
\begin{align}
  \frac{N_L^{\pi e}}
       {N_L^{\pi\mu}}
  &\sim10^{-4},
  \\
  \frac{(N_L^{\pi e})^2}
       {N_H^{\pi e}N_L^{\pi\mu}}
  &\sim10^{-8},
\end{align}
and are therefore negligible at the required precision. Keeping only the
leading correction term gives
\begin{equation}
  R_{e/\mu}^{\mathrm{true}}
  \simeq
  \frac{N_H}{N_L}
  \left(
    1+
    \frac{N_L^{\pi e}}
         {N_H^{\pi e}}
  \right).
\end{equation}

Defining the tail correction
\begin{equation}
  C_T
  \equiv
  \frac{N_L^{\pi e}}
       {N_H^{\pi e}},
\end{equation}
the final expression becomes
\begin{equation}
  R_{e/\mu}^{\mathrm{true}}
  \simeq
  \frac{N_H}{N_L}
  \left(
    1+C_T
  \right),
  \label{eq:tail-corrected-remu}
\end{equation}
before relative-acceptance effects.


\subsection{Relative Efficiency and Acceptance}
\label{subsec:remu-efficiency-correction}

Let \(\varepsilon_H\) be the probability that a true \(\pi\to e\nu\) event is
accepted, reconstructed, and counted in \(N_H\). Let \(\varepsilon_L\) be the
corresponding probability for a true \(\pi\to\mu\to e\) event counted in
\(N_L\). The relative efficiency correction is
\begin{equation}
  R^\varepsilon
  \equiv
  \frac{\varepsilon_H}{\varepsilon_L}.
  \label{eq:relative-efficiency-definition}
\end{equation}
The corrected estimator is therefore
\begin{equation}
  R_{e/\mu}
  \simeq
  \frac{N_H}{N_L}
  \left(1+C_T\right)
  R^\varepsilon .
  \label{eq:remu-corrected-estimator-appendix}
\end{equation}

\subsection{Angular Dependence}
\label{subsec:remu-angular-dependence}

The tail and efficiency corrections should be understood as functions of the
positron direction. Although the detector geometry is approximately symmetric
under rotations about the beam axis, the response is not fully uniform because
of the forward beam entrance region, upstream material, calorimeter
containment, reconstruction thresholds, and the directional segmentation of
the active target planes. As a result, many quantities entering the
\(R_{e/\mu}\) analysis depend on the reconstructed positron polar angle
\(\theta\).

A schematic angular-dependent form of the measurement is therefore
\begin{equation}
  R_{e/\mu}(\theta)
  \simeq
  \frac{N_H(\theta)}
       {N_L(\theta)}
  \times
  R^\varepsilon(\theta)
  \times
  \left(
    1+C_T(\theta)
  \right),
  \label{eq:angular-remu}
\end{equation}
where:
\begin{itemize}
  \item \(N_H(\theta)\) and \(N_L(\theta)\) are the reconstructed high- and
  low-energy event counts,
  \item \(C_T(\theta)\) is the low-energy
  \(\pi\to e\nu\) tail correction,
  \item \(R^\varepsilon(\theta)\) represents the relative efficiency and
  acceptance correction between the direct
  \(\pi\to e\nu\) and sequential
  \(\pi\to\mu\to e\) channels.
\end{itemize}

In practice, the analysis is performed in finite angular bins,
\begin{equation}
  R_{e/\mu}
  \simeq
  \sum_i
  w_i
  \times
  \frac{N_H^{(i)}}
       {N_L^{(i)}}
  \times
  R^\varepsilon_{(i)}
  \times
  \left(
    1+C_T^{(i)}
  \right),
  \label{eq:angular-binned-remu}
\end{equation}
where \(i\) labels the angular bin and the weights satisfy
\begin{equation}
  \sum_i w_i = 1 .
\end{equation}

This form emphasizes that the correction program is not purely an
energy-response problem. The measured ratio also depends on geometric
acceptance, reconstruction efficiency, and topology-dependent detector
response, all of which vary with positron direction. Once these corrections
are determined, their uncertainties and correlations must be propagated into
the final \(R_{e/\mu}\) result.

